\section{Реестр утверждений, происхождение свидетельств и открытые слоты}
\label{app:registers}

Канонический \src{observer/hypotheses.tsv}{реестр гипотез}
содержит 23 строки. Для каждой указаны идентификатор, зависимости, точная
формулировка, предпосылки, статус, адрес свидетельства и условие пересмотра.
Зависимости помещены до зависимых от них утверждений. Границы допустимых
историй являются частью формулировки. Например, опровержение переноса
исходного 11-битного примера на все слова совместимо с точностью на
выделенной четвёрке.

\begin{center}
\begin{tabularx}{\linewidth}{@{}l r X@{}}
\toprule
Статус в реестре & Из 23 & Адреса \\
\midrule
Доказано & 12 & OBS.01--03, R110.01, R110.05--09, R110.12, LOOP.01, E.01 \\
Опровергнуто & 4 & R110.02--04, R110.10 \\
Подтверждено в пределах опыта & 3 & AI.01--03 \\
Открыто & 2 & R110.11, AI.04 \\
\texttt{DECLARED\_NOT\_DERIVED} & 1 & E.02 \\
\texttt{UN} & 1 & E.03 \\
\bottomrule
\end{tabularx}
\end{center}

Статусы «доказано» и «опровергнуто» относятся к заявлению со всеми
предпосылками в соответствующей строке. У разных строк различны формы
свидетельства: математическое доказательство, вычислительно проверяемая
конечная конструкция или результат полного перебора. AI.01--03 фиксируют
соответственно связь ответа ассистента с показанным сигналом, расхождение
живого и записанного канала на новых запросах и пересмотр прежней стратегии
в проверенном семействе задач. Эти строки происходят из отдельных
испытаний ассистента. Счётчики замороженной OLMo в разделе~\ref{sec:olmo}
имеют собственные отчёты и панели. В E.01 доказана недоопределённость
добавленного предиката. E.02 хранит дополнительную посылку, E.03 ---
открытый статус конкретного экземпляра.

\paragraph{36 адресов утверждения.}
\src{observer/leaf/leaf.tsv}{Таблица слотов} сопоставляет шесть доменов
S (источник), I (идентичность), M (смысл), T (переход), A (действие),
O (наблюдение) с шестью позициями описания: происхождение, состояние,
носитель, локальное утверждение, свидетельство и предел. Для адресов этой
таблицы используем префикс \texttt{leaf:}. Например, \src{observer/leaf/leaf.tsv\#L25}{\texttt{leaf:T6}}
обозначает строку \texttt{law=T6} в TSV: шестая позиция домена переходов
отсылает к критерию для шумного канала (T3). Обозначение \hyperref[thm:conditional-bridge]{T6} без префикса
относится к условной теореме раздела~\ref{sec:boundary}. В срезе
32/36 строк имеют состояние \texttt{KNOWN}; у них заданы SHA-256 файла и
проверяемый адрес определения, функции или раздела текста.
Оставшиеся 4/36 имеют состояние \texttt{UNKNOWN} и пустые поля
источника-свидетельства. Их содержательная формулировка сохранена:
\begin{itemize}
\item \texttt{leaf:I6} --- численное тождество субъекта при ветвлении причинной истории.
Соглашение об учёте потомков задаёт идентификаторы ветвей.
\item \texttt{leaf:M6} --- независимое обоснование психофизической связи, принятой
как достаточное условие.
\item \texttt{leaf:A6} --- реализация всего бокового канала средствами одной
строго бинарной однородной Rule~110 для заявленного свидетеля (R110.11).
\item \texttt{leaf:O6} --- состояние феноменального предиката у данного экземпляра
в момент обратной связи (E.03).
\end{itemize}
\src{observer/leaf/render_leaf.py}{Валидатор изображения} проверяет порядок
адресов, соответствие хэшей источникам, существование символов, каноническое
представление TSV и его побайтное восстановление из SVG. Поэтому изображение
пригодно как проверяемая проекция таблицы.

\paragraph{Происхождение и статус релиза.}
В \src{SHA256SUMS}{корневом манифесте выпуска} приведены контрольные суммы PDF, обоих журналов Lean и исходников. Для корпуса наблюдателя
\src{observer/sha256.txt}{отдельный манифест} фиксирует 48 файлов. Для исходной редакции отдельно сохранён
\src{observer/provenance/observer-original.sha256}{исторический список
контрольных сумм} из 36 записей. В подготовке публикационного комплекта
использованы 35 исходных файлов этого списка. Запись служебного
исследовательского задания сохранена в списке для учёта происхождения;
само задание в публичный комплект не включено. Изменения навигации и
адресов свидетелей отражены в публикационных хэшах. Сохранённые
JSON-отчёты содержат исходные наблюдения. Контракт, использованный в
прежних запусках, проверяется по байтам его текста из отчётной редакции.

\paragraph{Условия продолжения.}
По R110.11 нужна проверяемая кодировка ввода, хранения, доставки и
ответа внутри одной бинарной динамики с траекториями всех четырёх историй.
По AI.04 нужны новая задача и проверка переноса внутреннего состояния при
закрытых исходных позициях, целой истории в положительном контроле,
обратимой подмене и стирании. Контракт численного G4 приведён в
приложении~\ref{app:gates}. Чтобы локализовать перенос в выбранном
подпространстве, до опыта дополнительно задаются слой, вмешательство, мера
посредничества, контролируемые альтернативы и допуск. Результат для одного
текстового носителя этого условия не закрывает.

Полная проверка канонических импортов наблюдателя выполнена на Lean~4.35.0
при подготовке монографии (приложение~\ref{app:lean}). Статусы панелей OLMo
относятся к экспериментальному dev-срезу. Для гипотезы устойчивой
пластичности требуется межэпизодное
изменение организации носителя и подтверждённое причинное действие этого
изменения на последующий вход при контроле внешней истории. Для E.02
нужны независимые основания выбора моста либо различающиеся операциональные
предсказания конкурирующих мостов с явными связывающими посылками.
Статус E.03 остаётся \texttt{UN} в указанном срезе.
